\MakeAbstract{
    建筑幕墙是现代建筑的重要外围护结构，其服役状态直接影响建筑安全和运维质量。现实工程中，幕墙巡检长期依赖人工经验，存在效率低、主观性强、难以处理海量图像等问题。已有计算机视觉方法能够完成裂缝、锈蚀、破损等局部缺陷检测，但多数方案仍停留在单模型、单图像、单缺陷识别层面，只能检测某张图像是否存在某类缺陷，缺少多组巡检图像、多种幕墙材质、多类检测结果之间的汇总融合。对于土木工程领域用户而言，实际需求并不是孤立的单图检测结论，而是面向整栋建筑形成的多维度韧性评估结论。因此，为打通从项目级图像资产管理、单图检测结果到建筑级韧性评估报告的完整链路，本文设计并实现了一套基于 AI Agent 的建筑幕墙韧性评估软件系统。

    本文的研究重点是 AI Agent 参与工程任务规划、能力调度与信息整合的系统性方法。本文工作的主要亮点包括两个方面：其一，提出面向建筑幕墙多材质图像的 Agent 自主任务调度方法。该方法由分类 Agent 根据图像内容自主判断需要执行的原子检测能力，并将开放语义判断约束为工程系统可消费的任务代码；其二，提出 LLM 分层信息整合策略。该策略将所有多源检测结果、项目背景、图像分组和人工复核意见等信息进行分层整合，将大语言模型从简单文本生成工具扩展为“路由、压缩、汇总”的智能节点，有效解决了海量数据直接输入 LLM 模型带来的上下文过载和重点淹没问题。

    在工程实现上，原型系统主要使用 Golang、Python 和 TypeScript 编程语言、React、Gin 和 PyTorch 框架、gRPC 服务通信协议、RocketMQ 消息队列、WebSocket 实时消息推送、MySQL 关系型数据库、Redis 非关系型数据库、MinIO 对象存储和 Docker 容器化等技术。工程链路为 Agent 提供高可用、高性能、高并发的运行环境，原型系统提供从建筑图像资产构建、Agent 智能检测、人工复核到韧性评估报告的完整闭环工作链。作者现已完成原型系统的部署和实际场景试用，获得同济大学土木工程学院科研团队认可，为建筑幕墙数字运维提供了一种兼具工程可实施性和智能决策能力的解决方案。
}{AI Agent，建筑幕墙，任务编排，韧性评估}

\MakeAbstractEng{
    Building curtain walls are important exterior envelope structures in modern buildings, and their service condition directly affects building safety and operation and maintenance quality. In real engineering practice, curtain wall inspection has long relied on manual experience, which leads to low efficiency, strong subjectivity, and difficulty in processing massive image data. Existing computer vision methods can detect local defects such as cracks, corrosion, and breakage. However, most of them still remain at the level of single-model, single-image, and single-defect recognition. They can only detect whether a certain defect exists in a certain image, and lack the ability to summarize and integrate multiple groups of inspection images, different curtain wall materials, and multiple types of detection results. For civil engineering users, the actual requirement is not an isolated single-image detection conclusion, but a multi-dimensional resilience assessment conclusion for the whole building. Therefore, to bridge the complete workflow from project-level image asset management and single-image detection results to building-level resilience assessment report generation, this thesis designs and implements an AI Agent-based software system for building curtain wall resilience assessment.

    The research focus of this thesis is the systematic method by which AI Agents participate in engineering task planning, capability scheduling, and information integration. The main highlights of this work include two aspects. First, an Agent autonomous task scheduling method is proposed for multi-material building curtain wall images. In this method, a classification Agent autonomously determines the atomic detection capabilities that need to be executed according to image content, and constrains open-ended semantic judgments into task codes that can be consumed by the engineering system. Second, an LLM hierarchical information integration strategy is proposed. This strategy hierarchically integrates multi-source detection results, project background, image grouping, human review comments, and other information. It extends large language models from simple text generation tools into intelligent nodes for routing, compression, and aggregation, thereby effectively solving the context overload and key-information dilution caused by directly feeding massive data into LLMs.

    In terms of engineering implementation, the prototype system mainly uses Golang, Python, and TypeScript as programming languages, React, Gin, and PyTorch as frameworks, gRPC for service communication, RocketMQ as the message queue, WebSocket for real-time message pushing, MySQL as relational database storage, Redis as non-relational database storage, MinIO as object storage, and Docker for containerized deployment. The engineering workflow provides a highly available, high-performance, and highly concurrent runtime environment for Agents. The prototype system provides a complete closed-loop workflow from building image asset construction, Agent intelligent detection, and human review to resilience assessment report generation. The prototype system has been deployed and tested in practical scenarios, and has received recognition from the research team of the College of Civil Engineering at Tongji University. This work provides a solution for digital operation and maintenance of building curtain walls that combines engineering feasibility with intelligent decision-making capability.
}{AI Agent, Building Curtain Wall, Task Orchestration, Resilience Assessment}
